% SOLO P2.
% BORRADOR redactado por Alessandra a partir del protocolo de P02 (Marcia Arce) para que
% el informe compile completo. Describe lo que hacen P02 y src/protocolo.py; Marcia puede
% reemplazarlo o ajustarlo con su propia redacción.
\section{Diseño experimental y evaluación}
\label{sec:diseno}

\subsection{Particiones}
Los 7043 clientes se dividieron una sola vez, de forma estratificada por \var{Churn} y con semilla 42,
en entrenamiento (60\%, 4225 clientes), validación (20\%, 1409) y prueba (20\%, 1409). La
estratificación mantiene la proporción de abandono de 26.5\% en los tres conjuntos (1121, 374 y 374
abandonos). Cada conjunto cumple un único papel: con entrenamiento se ajustan los preprocesadores y
los modelos; con validación se comparan las configuraciones y se elige una por modelo; prueba se
reserva para una sola evaluación final de los tres modelos ya elegidos.

\subsection{Búsqueda de hiperparámetros}
Se recorrieron de forma exhaustiva las grillas de la Tabla~\ref{tab:grilla}. Cada configuración se
ajustó con entrenamiento y se evaluó en validación. Se usó una validación fija (\emph{hold-out}) en
lugar de validación cruzada de $k$ pliegues. Esta decisión simplifica la coordinación entre los
integrantes y permite reproducir exactamente las mismas particiones, aunque hace que los resultados
dependan de la partición elegida (Sección~\ref{sec:discusion}).

El criterio de selección, fijado antes de entrenar, fue el mayor F1-score de \var{Churn=1} en
validación. Ante un empate se prefiere la configuración más simple: menor $C$ en la regresión
logística; menor profundidad y, luego, hojas más grandes en el árbol. Los modelos finales conservan
el ajuste hecho solo con entrenamiento, de modo que lo que se evalúa en prueba es exactamente lo que
se seleccionó.

\subsection{Métricas}
La clase positiva es \var{Churn=1} y la etiqueta se asigna con el umbral 0.5 sobre la probabilidad
estimada. A partir de los verdaderos positivos (VP), falsos positivos (FP), falsos negativos (FN) y
verdaderos negativos (VN) de la matriz de confusión se reportan:
\begin{itemize}
    \item \emph{Accuracy}: proporción total de aciertos.
    \item \emph{Precision}, $\mathrm{VP}/(\mathrm{VP}+\mathrm{FP})$: de los clientes que el modelo
    señala, cuántos abandonan.
    \item \emph{Recall}, $\mathrm{VP}/(\mathrm{VP}+\mathrm{FN})$: de los clientes que abandonan,
    cuántos señala el modelo.
    \item F1-score: media armónica de \emph{Precision} y \emph{Recall}.
    \item ROC-AUC: probabilidad de que un cliente que abandona reciba una probabilidad estimada mayor
    que uno que permanece. No depende del umbral.
\end{itemize}

La exactitud por sí sola es engañosa con clases desbalanceadas: el clasificador de referencia alcanza
0.735 sin detectar a nadie. Además, en retención los dos errores no cuestan lo mismo. Un falso
negativo es un cliente que se pierde sin haber sido contactado; un falso positivo es una oferta de
retención innecesaria. Como esos costos no se conocen, se usó F1 como criterio de selección porque
exige un equilibrio entre ambos errores.

\subsection{Incertidumbre e importancia de variables}
Con 374 abandonos en prueba, las métricas tienen un error de muestreo apreciable. Se estimó con
\emph{bootstrap}~\cite{bootstrap}: 2000 remuestreos con reemplazo de los clientes de prueba e
intervalos percentil del 95\%. Para comparar dos modelos, la diferencia se calcula en cada remuestreo
sobre los mismos clientes (bootstrap pareado) y se considera distinguible solo si su intervalo
excluye el cero.

Para comparar qué variables usa cada modelo en una misma escala se calculó la importancia por
permutación~\cite{rf}: se desordena una de las 19 variables originales en validación y se mide cuánto
cae el F1-score, en promedio sobre 30 repeticiones.
